% !TEX root = ../main.tex
\chapter{Contr\^ole de Version --- Git et DVC}
\label{ch:version-control}

\section{Introduction}

Dans tout projet de Machine Learning, la reproductibilit\'e est fondamentale.
Sans un syst\`eme de contr\^ole de version rigoureux, il devient impossible de
retracer quelle version du code, des donn\'ees et des hyperparam\`etres a produit
un mod\`ele donn\'e. Ce chapitre couvre les outils essentiels~: \textbf{Git} pour
le code et \textbf{DVC} (Data Version Control) pour les donn\'ees et les mod\`eles.

\begin{attention}[title=\textbf{Le pi\`ege du <<~c'\'etait mieux avant~>>}]
Sans versionnement, vous ne pourrez jamais revenir \`a un \'etat ant\'erieur de
votre projet. Un mod\`ele performant peut \^etre perdu \`a jamais si vous \'ecrasez
les fichiers sans pr\'ecaution.
\end{attention}

%% ============================================================
\section{Git~: les fondamentaux pour le ML}
\label{sec:git-essentials}
%% ============================================================

\subsection{Rappels sur Git}
\index{Git}

\begin{definition}[D\'ep\^ot Git]
Un \textbf{d\'ep\^ot Git} (\emph{repository}) est un r\'epertoire dont l'historique
complet des modifications est enregistr\'e. Chaque instantan\'e est appel\'e un
\textbf{commit}.
\end{definition}

Les commandes essentielles~:

\begin{codeblock}[title=\textbf{Commandes Git de base}]
\begin{minted}{bash}
git init                    # Initialiser un d\'ep\^ot
git add <fichier>           # Ajouter au staging
git commit -m "message"     # Cr\'eer un commit
git log --oneline --graph   # Visualiser l'historique
git diff                    # Voir les modifications
\end{minted}
\end{codeblock}

\subsection{Branches et fusion}
\index{Git!branches}

Les branches permettent de travailler sur des fonctionnalit\'es ou des
exp\'eriences en parall\`ele sans affecter la branche principale.

\begin{codeblock}[title=\textbf{Gestion des branches}]
\begin{minted}{bash}
git branch feature/new-model    # Cr\'eer une branche
git checkout feature/new-model  # Basculer sur la branche
git checkout -b feature/xgboost # Cr\'eer et basculer
git merge feature/new-model     # Fusionner dans la branche courante
git branch -d feature/new-model # Supprimer la branche
\end{minted}
\end{codeblock}

\subsection{Le fichier \file{.gitignore} pour le ML}
\index{.gitignore}

En ML, de nombreux fichiers ne doivent \textbf{jamais} \^etre commit\'es~:

\begin{codeblock}[title=\textbf{Exemple de \file{.gitignore} pour un projet ML}]
\begin{minted}{bash}
# Donn\'ees volumineuses
data/raw/
data/processed/
*.csv
*.parquet
*.h5

# Mod\`eles entra\^in\'es
models/*.pkl
models/*.pt
models/*.onnx

# Environnements
venv/
.conda/
__pycache__/

# Notebooks checkpoints
.ipynb_checkpoints/

# Secrets
.env
*.key
\end{minted}
\end{codeblock}

%% ============================================================
\section{Workflow Git pour les projets ML}
\label{sec:git-workflow-ml}
%% ============================================================

\index{Git!workflow ML}

\begin{bonnepratique}[title=\textbf{Strat\'egie de branches pour le ML}]
\begin{itemize}
  \item \texttt{main}~: code stable, mod\`eles valid\'es en production.
  \item \texttt{develop}~: int\'egration continue des nouvelles fonctionnalit\'es.
  \item \texttt{feature/*}~: nouvelles fonctionnalit\'es ou architectures.
  \item \texttt{experiment/*}~: exp\'eriences exploratoires (peuvent \^etre abandonn\'ees).
  \item \texttt{data/*}~: modifications aux pipelines de donn\'ees.
\end{itemize}
\end{bonnepratique}

\begin{bonnepratique}[title=\textbf{Revues de code (Pull Requests) en ML}]
Une PR en ML devrait inclure~:
\begin{enumerate}
  \item Les m\'etriques avant/apr\`es modification.
  \item Une justification de l'approche choisie.
  \item Les visualisations pertinentes (courbes d'apprentissage, matrices de confusion).
  \item La reproductibilit\'e~: graines al\'eatoires fix\'ees, environnement document\'e.
\end{enumerate}
\end{bonnepratique}

%% ============================================================
\section{DVC~: Data Version Control}
\label{sec:dvc}
%% ============================================================

\index{DVC}

\subsection{Pourquoi Git ne suffit pas}

\begin{attention}[title=\textbf{Limites de Git pour les fichiers volumineux}]
Git stocke l'int\'egralit\'e de chaque version d'un fichier. Pour un jeu de
donn\'ees de 10\,Go modifi\'e 5 fois, le d\'ep\^ot occuperait 50\,Go.
Les plateformes comme GitHub imposent une limite de 100\,Mo par fichier.
\end{attention}

\begin{definition}[DVC]
\textbf{DVC} (Data Version Control) est un outil open-source qui \'etend Git
pour g\'erer les fichiers volumineux (donn\'ees, mod\`eles). Il remplace les
fichiers par des pointeurs l\'egers (\file{.dvc}) et stocke le contenu
r\'eel sur un stockage distant (S3, GCS, Azure, SSH, etc.).
\end{definition}

\subsection{Installation et initialisation}

\begin{codeblock}[title=\textbf{Installation de DVC}]
\begin{minted}{bash}
pip install dvc[s3]      # Avec support S3
pip install dvc[gs]      # Avec support Google Cloud
pip install dvc[all]     # Tous les backends

# Initialiser DVC dans un d\'ep\^ot Git existant
cd mon-projet-ml
dvc init
git add .dvc .dvcignore
git commit -m "Initialiser DVC"
\end{minted}
\end{codeblock}

\subsection{Commandes fondamentales de DVC}
\index{DVC!commandes}

\begin{formules}[title=\textbf{Commandes DVC essentielles}]
\begin{tabular}{ll}
\toprule
\textbf{Commande} & \textbf{Description} \\
\midrule
\cli{dvc init} & Initialiser DVC dans un d\'ep\^ot Git \\
\cli{dvc add <fichier>} & Suivre un fichier/dossier avec DVC \\
\cli{dvc push} & Envoyer les donn\'ees vers le stockage distant \\
\cli{dvc pull} & R\'ecup\'erer les donn\'ees depuis le stockage distant \\
\cli{dvc remote add} & Configurer un stockage distant \\
\cli{dvc status} & V\'erifier l'\'etat des fichiers suivis \\
\cli{dvc checkout} & Restaurer les fichiers \`a partir du cache \\
\cli{dvc gc} & Nettoyer le cache local \\
\bottomrule
\end{tabular}
\end{formules}

\begin{codeblock}[title=\textbf{Utilisation de DVC au quotidien}]
\begin{minted}{bash}
# Configurer le stockage distant
dvc remote add -d myremote s3://mon-bucket/dvc-store
git add .dvc/config
git commit -m "Configurer le remote DVC"

# Suivre un jeu de donn\'ees
dvc add data/raw/dataset.csv
git add data/raw/dataset.csv.dvc data/raw/.gitignore
git commit -m "Ajouter dataset v1"

# Envoyer les donn\'ees
dvc push

# R\'ecup\'erer les donn\'ees (apr\`es un clone)
git clone <url-du-repo>
dvc pull
\end{minted}
\end{codeblock}

\subsection{Pipelines DVC}
\index{DVC!pipelines}

DVC permet de d\'efinir des pipelines reproductibles via un fichier \file{dvc.yaml}.

\begin{codeblock}[title=\textbf{Exemple de \file{dvc.yaml}}]
\begin{minted}{yaml}
stages:
  prepare:
    cmd: python src/prepare.py
    deps:
      - src/prepare.py
      - data/raw/dataset.csv
    outs:
      - data/processed/train.csv
      - data/processed/test.csv

  train:
    cmd: python src/train.py
    deps:
      - src/train.py
      - data/processed/train.csv
    params:
      - train.learning_rate
      - train.n_estimators
    outs:
      - models/model.pkl
    metrics:
      - metrics/scores.json:
          cache: false

  evaluate:
    cmd: python src/evaluate.py
    deps:
      - src/evaluate.py
      - models/model.pkl
      - data/processed/test.csv
    metrics:
      - metrics/eval.json:
          cache: false
    plots:
      - metrics/confusion_matrix.csv:
          x: predicted
          y: actual
\end{minted}
\end{codeblock}

\begin{codeblock}[title=\textbf{Ex\'ecution et comparaison de pipelines}]
\begin{minted}{bash}
# Ex\'ecuter le pipeline complet
dvc repro

# Comparer les m\'etriques entre branches/commits
dvc metrics diff
dvc params diff

# Visualiser le DAG du pipeline
dvc dag
\end{minted}
\end{codeblock}

\subsection{Fichier \file{.dvcignore}}
\index{.dvcignore}

Comme \file{.gitignore}, le fichier \file{.dvcignore} indique \`a DVC les
fichiers \`a ignorer lors du suivi de r\'epertoires~:

\begin{codeblock}[title=\textbf{Exemple de \file{.dvcignore}}]
\begin{minted}{bash}
# Fichiers temporaires
*.tmp
*.log

# Caches
__pycache__/
.ipynb_checkpoints/

# Fichiers syst\`eme
.DS_Store
Thumbs.db
\end{minted}
\end{codeblock}

%% ============================================================
\section{Architecture Git + DVC}
\label{sec:git-dvc-architecture}
%% ============================================================

\begin{figure}[htbp]
\centering
\begin{tikzpicture}[
    node distance=1.5cm and 2.5cm,
    block/.style={rectangle, draw, rounded corners, minimum width=3cm,
                  minimum height=1cm, align=center, font=\small},
    storage/.style={cylinder, draw, shape border rotate=90,
                    minimum width=2cm, minimum height=1.5cm,
                    align=center, font=\small},
    arrow/.style={-{Stealth}, thick},
    label/.style={font=\footnotesize\itshape, text=gray}
  ]

  % D\'eveloppeur
  \node[block, fill=blue!15] (dev) {D\'eveloppeur\\local};

  % Git
  \node[block, fill=green!15, right=of dev] (git) {D\'ep\^ot Git\\(GitHub/GitLab)};
  \node[label, above=0.2cm of git] {Code + fichiers \file{.dvc}};

  % DVC Remote
  \node[storage, fill=orange!15, right=of git] (remote) {Stockage\\distant\\(S3/GCS)};
  \node[label, above=0.2cm of remote] {Donn\'ees + mod\`eles};

  % Cache local
  \node[storage, fill=yellow!15, below=of dev] (cache) {Cache\\DVC\\local};

  % Fl\`eches
  \draw[arrow] (dev) -- node[above, font=\scriptsize] {\cli{git push}} (git);
  \draw[arrow] (git) -- node[below, font=\scriptsize] {\cli{git pull}} (dev);

  \draw[arrow] (dev) -- node[left, font=\scriptsize] {\cli{dvc add}} (cache);
  \draw[arrow] (cache) -- node[right, font=\scriptsize] {\cli{dvc checkout}} (dev);

  \draw[arrow, bend left=20] (cache.east) to
    node[below, font=\scriptsize] {\cli{dvc push}} (remote.south);
  \draw[arrow, bend left=20] (remote.south) to
    node[above, font=\scriptsize] {\cli{dvc pull}} (cache.east);

\end{tikzpicture}
\caption{Architecture du workflow Git + DVC. Le code et les fichiers
\file{.dvc} transitent par Git, tandis que les donn\'ees volumineuses sont
g\'er\'ees par le cache DVC local et le stockage distant.}
\label{fig:git-dvc-workflow}
\end{figure}

%% ============================================================
\section{Comparaison des outils de versionnement de donn\'ees}
\label{sec:version-tools-comparison}
%% ============================================================

\index{Git LFS}\index{Delta Lake}\index{LakeFS}

\begin{table}[htbp]
\centering
\caption{Comparaison des outils de versionnement de donn\'ees}
\label{tab:version-tools}
\small
\begin{tabular}{lcccc}
\toprule
\textbf{Crit\`ere} & \textbf{DVC} & \textbf{Git LFS} & \textbf{Delta Lake} & \textbf{LakeFS} \\
\midrule
Licence & Open-source & Open-source & Open-source & Open-source \\
Stockage & S3, GCS, etc. & Serveur LFS & Object store & S3-compatible \\
Pipelines & Oui & Non & Non & Non \\
Branching donn\'ees & Via Git & Via Git & Non & Oui (natif) \\
Taille max fichier & Illimit\'ee & Illimit\'ee & Illimit\'ee & Illimit\'ee \\
Int\'egration Git & Excellente & Native & Aucune & Bonne \\
Courbe d'apprentissage & Mod\'er\'ee & Facile & \'Elev\'ee & Mod\'er\'ee \\
Cas d'usage id\'eal & Projets ML & Binaires & Data lakes & Data lakes \\
\bottomrule
\end{tabular}
\end{table}

\begin{remarque}
\textbf{Git LFS} est plus simple que DVC mais ne supporte pas les pipelines
de donn\'ees. \textbf{LakeFS} offre un v\'eritable branching au niveau des
donn\'ees, ce qui est id\'eal pour les lacs de donn\'ees \`a grande \'echelle.
\end{remarque}

%% ============================================================
\section{Mini-projet~: versionner un dataset et un mod\`ele}
\label{sec:miniprojet-dvc}
%% ============================================================

\begin{miniprojet}[title=\textbf{Versionner un projet ML complet avec Git et DVC}]
\textbf{Objectif}~: Mettre en place un workflow Git + DVC pour un projet de
classification.

\begin{enumerate}
  \item Initialiser un d\'ep\^ot Git et DVC~:
\begin{minted}{bash}
mkdir projet-classification && cd projet-classification
git init && dvc init
dvc remote add -d local_remote /tmp/dvc-storage
\end{minted}

  \item Cr\'eer la structure du projet~:
\begin{minted}{bash}
mkdir -p data/raw data/processed models src metrics
\end{minted}

  \item Ajouter un jeu de donn\'ees et le suivre avec DVC~:
\begin{minted}{python}
# src/download_data.py
from sklearn.datasets import load_wine
import pandas as pd

data = load_wine(as_frame=True)
df = pd.concat([data.data, data.target], axis=1)
df.to_csv("data/raw/wine.csv", index=False)
\end{minted}

\begin{minted}{bash}
python src/download_data.py
dvc add data/raw/wine.csv
git add data/raw/wine.csv.dvc data/raw/.gitignore
git commit -m "Ajouter wine dataset v1"
dvc push
\end{minted}

  \item Cr\'eer un pipeline \file{dvc.yaml} avec les \'etapes \texttt{prepare},
        \texttt{train} et \texttt{evaluate}.

  \item Ex\'ecuter le pipeline, modifier un hyperparam\`etre et comparer les
        r\'esultats~:
\begin{minted}{bash}
dvc repro
dvc metrics show
# Modifier params.yaml, puis :
dvc repro
dvc metrics diff
\end{minted}
\end{enumerate}
\end{miniprojet}

%% ============================================================
\section{Exercices}
\label{sec:exos-version-control}
%% ============================================================

\begin{exercice}[Cr\'eer un \file{.gitignore} ML]
\textbf{1/3}
Cr\'eez un fichier \file{.gitignore} adapt\'e \`a un projet de Deep Learning
utilisant PyTorch. Incluez les patterns pour les donn\'ees, les checkpoints,
les logs TensorBoard, les environnements virtuels et les fichiers syst\`eme.
\end{exercice}

\begin{exercice}[Pipeline DVC multi-\'etapes]
\textbf{2/3}
Cr\'eez un pipeline DVC complet pour un projet de NLP comprenant~:
\begin{enumerate}
  \item T\'el\'echargement des donn\'ees brutes.
  \item Pr\'etraitement (tokenisation, nettoyage).
  \item Entra\^inement d'un mod\`ele de classification de texte.
  \item \'Evaluation avec m\'etriques (accuracy, F1-score).
\end{enumerate}
Utilisez \file{params.yaml} pour les hyperparam\`etres et comparez deux
configurations avec \cli{dvc metrics diff}.
\end{exercice}

\begin{exercice}[Branching et exp\'eriences parall\`eles]
\textbf{3/3}
Mettez en place un workflow o\`u deux exp\'eriences sont men\'ees en parall\`ele
sur des branches Git diff\'erentes. Chaque branche utilise une architecture
de mod\`ele diff\'erente (par exemple, Random Forest vs.\ Gradient Boosting).
Utilisez DVC pour~:
\begin{itemize}
  \item Suivre les donn\'ees d'entra\^inement (partag\'ees entre branches).
  \item Suivre les mod\`eles entra\^in\'es (propres \`a chaque branche).
  \item Comparer les m\'etriques entre les deux branches avec
        \cli{dvc metrics diff}.
  \item Fusionner la meilleure branche dans \texttt{main}.
\end{itemize}
\end{exercice}

%% ============================================================
\section{Aide-m\'emoire}
\label{sec:cheatsheet-version-control}
%% ============================================================

\begin{cheatsheet}[title=\textbf{Git + DVC --- Aide-m\'emoire}]
\small
\begin{tabular}{p{4.5cm}p{8.5cm}}
\toprule
\textbf{Action} & \textbf{Commande} \\
\midrule
Initialiser Git + DVC & \cli{git init \&\& dvc init} \\
Configurer remote DVC & \cli{dvc remote add -d nom s3://bucket/path} \\
Suivre un fichier & \cli{dvc add data/fichier.csv} \\
Commiter les pointeurs & \cli{git add fichier.csv.dvc .gitignore} \\
Envoyer donn\'ees & \cli{dvc push} \\
R\'ecup\'erer donn\'ees & \cli{dvc pull} \\
D\'efinir un pipeline & \'Editer \file{dvc.yaml} \\
Ex\'ecuter le pipeline & \cli{dvc repro} \\
Voir les m\'etriques & \cli{dvc metrics show} \\
Comparer les m\'etriques & \cli{dvc metrics diff} \\
Visualiser le DAG & \cli{dvc dag} \\
Cr\'eer une branche & \cli{git checkout -b experiment/nom} \\
Revue de code & Cr\'eer une Pull Request sur GitHub/GitLab \\
\bottomrule
\end{tabular}
\end{cheatsheet}
